% Indicate the main file. Must go at the beginning of the file.
% !TEX root = ../main.tex

%----------------------------------------------------------------------------------------
% CHAPTER TEMPLATE
%----------------------------------------------------------------------------------------


\chapter{Outlook and Future Work}
\label{ChapterOutlook}

This thesis demonstrated the feasibility of combining large language models (LLMs) with classical natural language processing (NLP) techniques to analyze large news corpora in a temporally-aware manner. While the proposed system successfully addresses the core research objectives and enables complex analytical queries over millions of documents, several limitations and opportunities for further development remain. This chapter outlines the most relevant directions for future work, focusing on analytical depth, robustness, scalability, and evaluation.

\section{Extending Analytical Capabilities}

In the current prototype, several NLP components, such as sentiment analysis, topic modeling, framing detection, and event detection are implemented using mocked outputs. This design choice allowed the thesis to focus on system architecture, orchestration, and agentic reasoning behavior rather than on individual model performance.

Future work should replace these mocked components with fully implemented, production ready NLP models. Transformer-based models fine-tuned on such large news corpora could improve sentiment and stance detection, while dynamic topic modeling methods (e.g., temporal extensions of BERTopic or LDA) could provide more fine-grained insights into topic evolution. Incorporating uncertainty estimation and calibration would further strengthen interpretability, particularly for ambiguous or polarized media coverage.


\section{Document Retrieval, Selection, and Scalability}

The system currently combines lexical retrieval via OpenSearch with an LLM-based document selection step. While effective in reducing large result sets to manageable and relevant subsets, this approach introduces computational overhead and may not scale optimally to substantially larger corpora.

Future iterations could integrate hybrid retrieval strategies that combine lexical search with dense vector retrieval, enabling improved semantic matching earlier in the pipeline. Learned re-ranking models or clustering-based sampling could reduce reliance on LLM-based filtering while preserving analytical quality.

\section{Parallelism and System Architecture}

The current pipeline executes most processing steps sequentially, despite the fact that many operations like document summarization, tool execution, and analysis could be parallelized. Introducing parallel execution strategies could significantly reduce runtime and improve system responsiveness.

Future work may explore multi-agent architectures in which specialized LLM agents handle distinct tasks, such as retrieval planning, document selection, analysis, and synthesis. Asynchronous execution and distributed task scheduling would further improve throughput and enable more interactive exploration of large datasets.

\section{Evaluation and Validation}

Evaluating the quality and correctness of complex analytical narratives generated by LLMs is inherently challenging. Particularly in the context of open-ended queries over large text corpora, traditional metrics such as accuracy or F1-score are insufficient. This thesis primarily relied on qualitative assessments through document inspection and case studies to validate system outputs.

Future research should introduce more systematic evaluation protocols, including human expert assessments and comparisons against baseline approaches such as keyword-based analysis or traditional topic modeling pipelines. Methods for detecting hallucinations and verifying faithfulness to retrieved documents would further strengthen the empirical validity of the system.

\section{Improving Robustness of Query Feasibility Assessment}

The system currently performs a preliminary feasibility assessment to determine whether a user query can be answered using the available corpus. While this early stopping mechanism avoids unnecessary computation, it relies entirely on a single LLM judgment and may prematurely reject valid questions.

A promising direction for future work is to replace this strict binary gate with a two-stage feasibility validation process. In such a design, the system would perform a lightweight retrieval step even when feasibility is uncertain and then evaluate the retrieved documents for semantic relevance and topical alignment. This hybrid approach would allow the system to fail more gracefully, reducing false negatives while still explicitly stating when a question cannot be answered from the available data.

\section{Broader Applicability}

Although this thesis focused on a large-scale news corpus, the proposed architecture is largely domain-agnostic and could be adapted to other text collections, such as scientific literature, legal documents, or corporate filings. Extensions to multilingual corpora or real-time media monitoring would further broaden the applicability of the system and open new research directions at the intersection of information retrieval, NLP, and large language models.

\section{Concluding Remarks}

Overall, the presented system represents an initial step toward agentic, LLM-assisted exploration of large text corpora. By addressing the outlined limitations and pursuing the proposed extensions, future work can enhance analytical depth, robustness, scalability, and scientific rigor, contributing to more transparent and effective tools for large-scale textual analysis.